\subsection{Global gauge group and one-family packet}

\begin{theorem}[Global Standard Model theorem]
The determinant-preserving connected stabilizer of the minimal carrier split is
\[
G_{\SM}^{\mathrm{glob}}
=
S(U(3)\times U(2))
\cong
\frac{SU(3)\times SU(2)\times U(1)_Y}{\mathbb Z_6}.
\]
\end{theorem}

\begin{proof}
An automorphism preserving the split $E_3\oplus E_2$ acts by a pair
$(u_3,u_2)\in U(3)\times U(2)$. Requiring determinant preservation on the full
carrier gives
\[
\det(u_3)\det(u_2)=1,
\]
so the connected stabilizer is exactly $S(U(3)\times U(2))$.

Define
\[
\Phi:SU(3)\times SU(2)\times U(1)_Y \longrightarrow S(U(3)\times U(2)),
\qquad
\Phi(g_3,g_2,z)=\bigl(z^2 g_3,\ z^{-3} g_2\bigr).
\]
This is a well-defined group homomorphism because the exponents are integers, and
\[
\det(z^2 g_3)\det(z^{-3} g_2)=z^6 z^{-6}=1.
\]

To prove surjectivity, take $(u_3,u_2)\in S(U(3)\times U(2))$. Choose
$z\in U(1)$ with
\[
z^6=\det(u_3)=\det(u_2)^{-1}.
\]
Then
\[
g_3:=z^{-2}u_3\in SU(3),
\qquad
g_2:=z^{3}u_2\in SU(2),
\]
and therefore
\[
\Phi(g_3,g_2,z)=(u_3,u_2).
\]

If $\Phi(g_3,g_2,z)=(\mathbf 1_3,\mathbf 1_2)$, then
\[
g_3=z^{-2}\mathbf 1_3,
\qquad
g_2=z^{3}\mathbf 1_2.
\]
The conditions $g_3\in SU(3)$ and $g_2\in SU(2)$ imply $z^6=1$. Hence
\[
\ker\Phi
=
\left\{
\bigl(z^{-2}\mathbf 1_3,\ z^{3}\mathbf 1_2,\ z\bigr): z^6=1
\right\}
\cong \mathbb Z_6.
\]
Therefore
\[
S(U(3)\times U(2))
\cong
\frac{SU(3)\times SU(2)\times U(1)_Y}{\mathbb Z_6}.
\]
\end{proof}
This is an internal carrier theorem: it identifies the connected stabilizer of the minimal carrier split. The paper does not claim that the global physical gauge group has already been experimentally fixed at the same level of closure.

Matter is read from the positive half-spinor
\[
S^+=\Lambda^{\mathrm{even}}E,
\qquad
\dim S^+=16.
\]

\begin{theorem}[One carrier family]
The positive half-spinor of the minimal carrier decomposes as
\[
S^+
\cong
(1,1)_0\oplus(3,2)_{1/6}\oplus(\bar 3,1)_{-2/3}\oplus(\bar 3,1)_{1/3}\oplus(1,2)_{-1/2}\oplus(1,1)_1,
\]
that is, exactly one chiral Standard Model family with $\nu^c$. Equivalently, in the familiar unified packaging language,
\[
S^+\cong \mathbf 1\oplus \mathbf{10}\oplus \overline{\mathbf 5}.
\]
\end{theorem}

\begin{proof}
One has
\[
S^+=\Lambda^0E\oplus\Lambda^2E\oplus\Lambda^4E.
\]
The $\Lambda^0E$ sector is immediately
\[
\Lambda^0E\cong (1,1)_0,
\]
which is identified with $\nu^c$ in left-chiral notation. For the minimal carrier,
\[
\Lambda^2E
\cong
\Lambda^2E_3\oplus(E_3\otimes E_2)\oplus\Lambda^2E_2.
\]
Here
\[
\Lambda^2E_3\cong \bar 3,
\qquad
E_3\otimes E_2\cong (3,2),
\qquad
\Lambda^2E_2\cong 1,
\]
and the hypercharges are obtained by adding the carrier charges from
\[
Y=\mathrm{diag}\!\left(-\frac13,-\frac13,-\frac13,\frac12,\frac12\right).
\]
Therefore
\[
\Lambda^2E
\cong
(\bar 3,1)_{-2/3}\oplus(3,2)_{1/6}\oplus(1,1)_1.
\]
Because the full determinant is trivial in $S(U(3)\times U(2))$, one has $\Lambda^5E\cong 1$ and hence
\[
\Lambda^4E\cong E^\ast.
\]
Equivalently,
\[
\Lambda^4E
\cong
(\Lambda^3E_3\otimes E_2)\oplus(\Lambda^2E_3\otimes\Lambda^2E_2).
\]
Using $\Lambda^3E_3\cong 1$, $\Lambda^2E_2\cong 1$, and $\Lambda^2E_3\cong E_3^\ast\cong\bar 3$, one gets
\[
\Lambda^4E\cong (1,2)_{-1/2}\oplus(\bar 3,1)_{1/3}.
\]
Collecting the three even-degree sectors yields
\[
S^+
\cong
(1,1)_0\oplus(\bar 3,1)_{-2/3}\oplus(3,2)_{1/6}\oplus(1,1)_1\oplus(\bar 3,1)_{1/3}\oplus(1,2)_{-1/2},
\]
with total multiplicity $1+3+6+1+3+2=16$. The $\Lambda^2E$ sector is the standard left-chiral $\mathbf{10}$ of $SU(5)$, the $\Lambda^4E$ sector is the standard left-chiral $\overline{\mathbf 5}$, and the $\Lambda^0E$ sector is the singlet $\mathbf 1$.
\end{proof}

In left-chiral notation the one-family decomposition takes the explicit form shown in \Cref{tab:one-family}. This is the place where the manuscript must be explicit: the carrier theorem is meaningful only if the reader can see how the sixteen states are organized.

\begin{table}[t]
\centering
\footnotesize
\begin{tabularx}{\linewidth}{@{}>{\raggedright\arraybackslash}p{0.18\linewidth}>{\raggedright\arraybackslash}p{0.19\linewidth}>{\raggedleft\arraybackslash}p{0.13\linewidth}>{\raggedright\arraybackslash}p{0.20\linewidth}>{\raggedleft\arraybackslash}p{0.08\linewidth}@{}}
\toprule
Exterior component & $\SM$ representation & Hypercharge & Field content & Mult. \\
\midrule
$\Lambda^0E$ & $(1,1)_0$ & $0$ & $\nu^c$ & $1$ \\
$\Lambda^2E_3$ & $(\bar 3,1)_{-2/3}$ & $-2/3$ & $u^c$ & $3$ \\
$E_3\otimes E_2$ & $(3,2)_{1/6}$ & $1/6$ & $Q_L$ & $6$ \\
$\Lambda^2E_2$ & $(1,1)_1$ & $1$ & $e^c$ & $1$ \\
\makecell[l]{$\Lambda^3E_3$\\$\otimes E_2$} & $(1,2)_{-1/2}$ & $-1/2$ & $L_L$ & $2$ \\
\makecell[l]{$\Lambda^2E_3$\\$\otimes\Lambda^2E_2$} & $(\bar 3,1)_{1/3}$ & $1/3$ & $d^c$ & $3$ \\
\midrule
\multicolumn{4}{r}{Total multiplicity} & $16$ \\
\bottomrule
\end{tabularx}
\caption{One chiral family from $S^+=\Lambda^{\mathrm{even}}E$, written in left-chiral notation with charge-conjugate singlets.}
\label{tab:one-family}
\end{table}

\begin{proposition}[Exact parity-code realization of the family packet]
\label{prop:parity-code-family}
Identifying each wedge monomial in $\Lambda E$ with its occupation word
\[
x=(x_1,\dots,x_5)\in\{0,1\}^5,
\]
the positive half-spinor is exactly the even-parity subspace
\[
S^+
\cong
\mathrm{span}\Bigl\{|x\rangle:\sum_{i=1}^5 x_i\equiv 0 \pmod 2\Bigr\}.
\]
Equivalently,
\[
S^+
=
\bigl(\Lambda^{\mathrm{even}}E_3\otimes\Lambda^{\mathrm{even}}E_2\bigr)
\oplus
\bigl(\Lambda^{\mathrm{odd}}E_3\otimes\Lambda^{\mathrm{odd}}E_2\bigr),
\]
and each summand has dimension $8$.
\end{proposition}

\begin{proof}
Choose a basis of $E$ adapted to the split $E_3\oplus E_2$. Exterior monomials are then in bijection with occupation words $x\in\{0,1\}^5$, and the exterior degree is exactly the Hamming weight $\sum_i x_i$. Hence $\Lambda^{\mathrm{even}}E$ is the span of the even-parity words. Splitting the degree into the $E_3$ and $E_2$ contributions gives the direct-sum decomposition above. Since
\[
\dim \Lambda^{\mathrm{even}}E_3=1+\binom32=4,
\qquad
\dim \Lambda^{\mathrm{odd}}E_3=\binom31+\binom33=4,
\]
and
\[
\dim \Lambda^{\mathrm{even}}E_2=1+\binom22=2,
\qquad
\dim \Lambda^{\mathrm{odd}}E_2=\binom21=2,
\]
each tensor-product sector has dimension $8$.
\end{proof}

\begin{remark}[Exact parity lock, not a loose code analogy]
The family packet is therefore not merely code-like: at the level of occupation kinematics it is exactly the classical even-parity $[5,4,2]$ code on five bits. What the present paper still does \emph{not} claim is a dynamical error-correction theorem for realistic noise channels. The exact statement is the parity lock itself. In the maximally mixed state on the $16$ occupation words, eight states lie in the $(\mathrm{even},\mathrm{even})$ sector and eight in $(\mathrm{odd},\mathrm{odd})$, so the coarse parity observables of the $3$-side and $2$-side carry exactly one bit of mutual information.
\end{remark}

The family packet can also be read as a two-dimensional occupancy map, shown in \Cref{fig:family-occupancy-map}. This is where the one-family table, the even-parity code, the split weight enumerator, and the $X=6Y$ arithmetic all meet in one glance.

\begin{figure}[t]
\centering
\footnotesize
\begin{tikzpicture}[>=Latex, x=1.7cm, y=1.45cm]
\draw[->, thick] (-0.2,0) -- (3.4,0) node[below] {$n_3$};
\draw[->, thick] (0,-0.2) -- (0,2.35) node[left] {$n_2$};
\foreach \x in {0,1,2,3} {
    \draw[gray!30] (\x,0) -- (\x,2.05);
    \node[below] at (\x,-0.02) {\x};
}
\foreach \y in {0,1,2} {
    \draw[gray!30] (0,\y) -- (3.05,\y);
    \node[left] at (-0.02,\y) {\y};
}
\node[align=left, anchor=west] at (2.05,2.15) {$n_3+n_2$ even\\$\Longleftrightarrow S^+$ admissible};

\foreach \x/\y in {1/0,3/0,0/1,2/1,1/2,3/2} {
    \draw[red!65!black, line width=0.6pt] (\x-0.08,\y-0.08) -- (\x+0.08,\y+0.08);
    \draw[red!65!black, line width=0.6pt] (\x-0.08,\y+0.08) -- (\x+0.08,\y-0.08);
}

\filldraw[blue!55!black] (0,0) circle (3.8pt);
\node[text=white, font=\scriptsize] at (0,0) {1};
\node[anchor=west] at (0.14,0.05) {$\nu^c$};

\filldraw[blue!55!black] (2,0) circle (4.5pt);
\node[text=white, font=\scriptsize] at (2,0) {3};
\node[anchor=west] at (2.12,0.08) {$u^c$};

\filldraw[green!55!black] (1,1) circle (5.5pt);
\node[text=white, font=\scriptsize] at (1,1) {6};
\node[anchor=west] at (1.14,1.08) {$Q_L$};

\filldraw[orange!80!black] (0,2) circle (3.8pt);
\node[text=white, font=\scriptsize] at (0,2) {1};
\node[anchor=west] at (0.12,1.9) {$e^c$};

\filldraw[orange!80!black] (3,1) circle (4.3pt);
\node[text=white, font=\scriptsize] at (3,1) {2};
\node[anchor=west] at (3.12,1.03) {$L_L$};

\filldraw[orange!80!black] (2,2) circle (4.5pt);
\node[text=white, font=\scriptsize] at (2,2) {3};
\node[anchor=west] at (2.12,1.9) {$d^c$};

\node[draw=black!55, rounded corners, fill=gray!4, align=center] at (1.55,-0.72) {%
    multiplicity $=\binom{3}{n_3}\binom{2}{n_2}$\\
    scaled hypercharge $X=3n_2-2n_3=6Y$};
\end{tikzpicture}
\caption{Occupancy map of the family packet. The admissible sectors are exactly the even lattice points in $(n_3,n_2)$, with multiplicities $(1,3,6,1,2,3)$ and field labels matching \Cref{tab:one-family}. This is the visual meeting point of parity lock, split weight enumerator, and family hypercharge arithmetic.}
\label{fig:family-occupancy-map}
\end{figure}

